\documentclass[UTF8,fontset=fandol,11pt,a4paper]{ctexart}
\usepackage[a4paper,margin=1.75cm,headheight=15pt]{geometry}
\usepackage{amsmath,amssymb,mathtools,booktabs,tabularx,array,enumitem}
\usepackage[most]{tcolorbox}\usepackage{tikz}\usetikzlibrary{arrows.meta,positioning}
\usepackage{xcolor,fancyhdr,hyperref,microtype,lastpage}
\newcolumntype{Y}{>{\raggedright\arraybackslash}X}\newcolumntype{L}[1]{>{\raggedright\arraybackslash}p{#1}}\newcolumntype{C}[1]{>{\centering\arraybackslash}p{#1}}
\setlength{\parindent}{2em}\setlength{\parskip}{0.34em}\setlength{\emergencystretch}{3em}\renewcommand{\arraystretch}{1.2}\setlength{\tabcolsep}{3pt}
\setlist[itemize]{itemsep=0.18em,topsep=0.35em,leftmargin=2em}\setlist[enumerate]{itemsep=0.18em,topsep=0.35em,leftmargin=2em}
\definecolor{deep}{RGB}{38,58,72}\definecolor{soft}{RGB}{244,247,248}\definecolor{greenish}{RGB}{52,91,74}\definecolor{warnc}{RGB}{136,61,58}\definecolor{purple}{RGB}{84,72,112}\definecolor{rulegray}{RGB}{110,116,120}
\hypersetup{colorlinks=true,linkcolor=deep,urlcolor=purple,pdftitle={CO-06 Cache 与存储层次}}
\pagestyle{fancy}\fancyhf{}\lhead{\small\color{deep}\textbf{408 · 计算机组成原理 CO-06}}\rhead{\small\color{deep}Locality · Identity · State · Cost}\cfoot{\small 第 \thepage\ 页 / 共 \pageref{LastPage} 页}\renewcommand{\headrulewidth}{0.4pt}
\newtcolorbox{corebox}[1]{breakable,colback=soft,colframe=deep,title=\textbf{#1},boxrule=0.8pt,arc=2mm}\newtcolorbox{methodbox}[1]{breakable,colback=white,colframe=greenish,title=\textbf{#1},boxrule=0.7pt,arc=1.5mm}\newtcolorbox{warnbox}[1]{breakable,colback=white,colframe=warnc,title=\textbf{#1},boxrule=0.7pt,arc=1.5mm}\newtcolorbox{examplebox}[1]{breakable,colback=white,colframe=purple,title=\textbf{#1},boxrule=0.7pt,arc=1.5mm}
\newcommand{\kw}[1]{\textbf{\color{deep}#1}}\newcommand{\code}[1]{\texttt{#1}}
\tikzset{co/.style={draw=deep,rounded corners=1.5mm,text width=2.7cm,minimum height=0.9cm,align=center,fill=soft,font=\small,inner sep=3pt},state/.style={co,double,double distance=1.2pt,fill=white},flow/.style={-{Latex[length=2mm]},thick,draw=deep},ref/.style={-{Latex[length=2mm]},dashed,draw=rulegray}}
\title{\vspace{-1.1cm}\textbf{Cache 与存储层次}\\[0.4em]\large 怎样维护一个正确的高速副本}
\author{408 · 计算机组成原理 Topic CO-06}\date{}
\begin{document}\maketitle
\begin{corebox}{母问题：小而快的副本怎样证明“它就是那个数据”？}
Cache 不是小型主存，而是带身份、有效性、修改状态和替换策略的副本系统：
\[
\boxed{\text{Locality}\to\text{Block Transfer}\to\text{Placement}\to\text{Tag / Valid}\to\text{Offset}\to\text{Replacement / Write State}\to\text{AMAT}}
\]
每次访问既要找到候选位置，又要证明候选行保存的是请求块，并决定副本修改何时向下传播。
\end{corebox}
\begin{methodbox}{本册定位与 Owner 边界}
\begin{itemize}
\item \kw{Owns：}局部性、块/行、tag/index/offset、直接/组相联/全相联、hit/miss 状态机、3C、替换、写策略、容量位数、AMAT 与层次接口。
\item \kw{Uses：}CO-05 的存储粒度与延迟；CO-07 的地址翻译接口；CO-04 的 miss stall/重叠。
\item \kw{Stops：}OS Page Cache、文件缓存、页表/置换和总线 DMA 不在本册重新定义。
\end{itemize}\end{methodbox}
\tableofcontents\newpage

\section{为什么“按地址直接找数据”不够}
Cache 容量小于下一级，多个主存块必须复用有限行。因此访问必须依次回答：允许放在哪？当前行属于哪个块？该块是否有效？块内取哪一字节？若被修改，谁拥有最新副本？这五问生成 placement、tag、valid、offset、dirty 和 replacement。

局部性是收益来源而非正确性证明：时间局部性预测近期重访，空间局部性预测邻近访问；预测失败仍必须通过 tag/valid 维持正确性。

\section{地址三分与组织}
设地址宽度为 $m$ 位，数据容量为 $C$ 字节，块大小为 $B$ 字节，相联度为 $E$，则组数
\[
S=\frac{C}{B\times E},\quad b=\log_2B,\quad s=\log_2S,\quad t=m-s-b.
\]
地址字段为 `tag | set index | block offset`。若按字编址或 $B$ 用字表示，先统一单位。

\begin{center}\small\begin{tabularx}{\linewidth}{L{2.1cm}C{1.5cm}Y Y}\toprule
组织 & 候选位置 & 优点 & 代价\\\midrule
直接映射 & 1 行 & 命中路径短、无需替换选择 & 冲突风险高\\
$E$ 路组相联 & 同组 $E$ 行 & 成本与冲突折中 & $E$ 个 tag 比较器和选择器\\
全相联 & 任意行 & 映射冲突最少 & 全部 tag 并行比较、替换成本高\\
\bottomrule\end{tabularx}\end{center}

\section{Hit/Miss 生命周期}
\subsection{Read hit}
索引定位候选组，比较 `valid=1` 且 tag 匹配的行，再由 offset 选择块内字节/字。命中只证明副本身份和有效性，不表示下一级没有其他副本。

\subsection{Read miss}
\[
\text{choose invalid/victim}\to\text{write back if dirty}\to\text{fetch block}\to\text{fill tag/data/state}\to\text{retry request}.
\]
顺序流水线可能 stall；能否继续执行由 CO-04 和实现的非阻塞结构决定。

\subsection{3C 模型}
compulsory 是第一次访问；capacity 是理想全相联同容量也装不下；conflict 是有限映射额外造成的互逐。分类必须绑定块大小、容量、相联度和访问序列，不能把所有 miss 都叫“Cache 太小”。

\section{写策略是两条独立轴}
Hit 时：write-through 同时更新下一级，流量大但一致关系直接；write-back 只更新 Cache 并设 dirty，逐出时才写回，流量较小但状态复杂。Miss 时：write-allocate 先取块再写入；no-write-allocate 绕过 Cache 写下一级。write-back 不自动推出 write-allocate。

\section{替换与状态成本}
直接映射没有选择；组相联/全相联需在候选中选 victim，可用 FIFO、Random、近似/精确 LRU。精确 LRU 的元数据和更新逻辑随相联度增长，不能笼统断言“每行只需 $\log_2E$ 位”。增大相联度降低冲突，可能增加 hit time、功耗和比较器数量。

\section{贯穿母例：两个块反复驱逐}
对直接映射 Cache，将地址除以 $B$ 得 block number，再算 $set=block\bmod S$。若两个 block 的 set 相同而 tag 不同，交替访问会互相驱逐，即使其他组为空。提高相联度让它们在同组共存，但需要额外比较与 victim 选择。

\begin{warnbox}{典型错误路径}
只看 tag 相同就判 hit，漏了 valid；只提高 Cache 容量就声称冲突一定消失；只看 hit rate 就宣称更快，漏了更长 hit time 或更大 miss penalty。
\end{warnbox}

\section{AMAT 与物理位数}
单级近似：
\[
AMAT=T_{hit}+MR\times MP.
\]
多级时 miss penalty 可能包含下一级 AMAT、写回脏块和总线占用。使用公式必须声明 hit time 是否包含并行 tag/data、MP 是额外时间、写回成本是否计入、stall 是否可重叠。

每行物理位数至少为 data、tag、valid，write-back 再加 dirty，替换/ECC 元数据按题设另算：
\[
C_{physical}=N_{lines}(data+tag+valid+dirty+metadata).
\]
“Cache 容量”若只指数据区，不能直接拿物理总位数替代。

\section{地址翻译接口}
PIPT、VIPT、VIVT 是 Cache 使用 VA/PA 与 TLB 的组织选择。VIPT 能并行的关键是索引位不依赖被翻译改变的 VPN 位；具体约束由页内偏移、块大小和相联度推出，不背固定上限。Cache hit 也不等于 TLB hit。

\section{五列概念边界}
\begin{center}\small\begin{tabularx}{\linewidth}{L{1.7cm}C{0.35cm}L{1.7cm}Y Y}\toprule
概念 A & $\neq$ & 概念 B & 真正区别与题面信号 & 混淆后果\\\midrule
Cache & $\neq$ & 主存 & 小容量带身份副本 vs 原始层级存储 & 把 tag/valid 当普通地址位\\
Tag & $\neq$ & Index & 身份证明 vs 候选组定位 & 地址三分错位\\
Valid & $\neq$ & Dirty & 行内容是否可用 vs 是否比下一级更新 & miss/写回条件错\\
Write-through & $\neq$ & Write-allocate & hit 更新传播轴 vs miss 是否取块轴 & 从一个策略推另一个\\
Hit rate & $\neq$ & Total time & 命中比例 vs hit/miss 成本综合 & 只按命中率判快慢\\
Cache miss & $\neq$ & TLB miss & 数据副本未命中 vs 翻译副本未命中 & 把重试处理者写错\\
\bottomrule\end{tabularx}\end{center}

\section{做题控制协议}
\begin{enumerate}
\item 统一编址单位、主存地址宽度、数据容量、块大小和相联度；
\item 算 line/set/offset/index/tag，并把访问转为 `(block,set,tag,offset)`；
\item 逐次推进 valid/tag/dirty/replacement 状态；
\item miss 写 victim、write-back、fill、策略和 retry；
\item AMAT 先声明成本模型，避免重复计算同一 penalty；
\item 用边界地址、冲突地址和替换后重访独立验证。
\end{enumerate}
\begin{methodbox}{压缩成第一动作}\
看到 Cache 题，先问：\kw{index 找到哪些候选？valid+tag 如何证明身份？offset 取哪一段？miss 时谁被替换、谁拥有最新数据？}
\end{methodbox}

\section{全科坐标与跨册交接}
\begin{corebox}{Map Coordinate：Data Location / Fast-Slow Path}
CO-06 负责回答“目标块现在在哪个副本、怎样证明身份、miss 后向哪一级请求”。它消费 CO-05 的存储服务模型，向 CO-07/B02 提供物理地址身份接口，并向 CO-04/I01 输出 hit/miss、fill、stall 与数据 ready 时刻。
\end{corebox}
Cache hit rate 不是完整 Cost；必须同时考虑 hit time、miss penalty、写回/填充流量、替换元数据和与流水线的重叠。它也不等于权限成功，更不等于 OS Page Cache 命中。

\section{考纲映射与来源处理}
\begin{center}\small\begin{tabularx}{\linewidth}{L{3.1cm}YY}\toprule
考点 & 本册机制 & 接口\\\midrule
局部性/层次 & block transfer、hit/miss、AMAT & CO-05\\
映射与替换 & 地址三分、3C、victim state & Rules\\
写策略/容量 & dirty、写回、物理位数 & CO-05/CO-04\\
TLB/Cache 组合 & PIPT/VIPT/VIVT 最小边界 & CO-07/CO-B02\\
\bottomrule\end{tabularx}\end{center}
\begin{corebox}{一句话}
Cache 的正确性来自“候选位置 + valid/tag 身份证明 + offset 定位 + 写状态管理”，性能才是在这个正确副本系统上由局部性、替换和 AMAT 产生的结果。\end{corebox}
本册吸收归档《Cache》《存储器层次》的地址、状态、替换、写策略和性能主干；固定 hit 周期、精确 LRU 位数、所有 miss 都阻塞整机等实现特例保留为题设条件。
\end{document}
